% !PRINTED 541
\weberpage{541}
%<<CONT
darstellbar sein, worin die $h_1, h_2, \ldots, h_n$ ganze rationale Zahlen
sind, und wenn diese Bedingung erfüllt ist, so ist auch umgekehrt $\eta$
mit $\xi$ nach dem Modul $\mu$ congruent. Ordnen wir die rechte Seite
mittelst (4) nach $\omega_1, \omega_2, \ldots, \omega_n$, so ergiebt sich
aus (5) und (6):
\weq{8}{\begin{array}{rclllll}
y_1 &=& x_1 &+ h_1a_{1,1} &+ h_2a_{1,2} &+ \cdots &+ h_na_{1,n} \\
y_2 &=& x_2 & &+ h_2a_{2,2} &+ \cdots &+ h_na_{2,n} \\
\multicolumn{7}{c}{.\quad.\quad.\quad.\quad.\quad.\quad.\quad.\quad.\quad.\quad.\quad.\quad.\quad.\quad.} \\
y_n &=& x_n & & & &+ h_na_{n,n}.
\end{array}}

Darin aber liegt der Beweis unserer Behauptung. Denn sind $y_1, y_2,
\ldots, y_n$ irgend welche gegebene ganze rationale Zahlen, so ist
zunächst aus
\[
y_n = x_n + h_na_{n,n}
\]
$x_n$ und $h_n$ vollständig bestimmt. Darauf erhalten wir aus der
vorletzten Gleichung
\[
y_{n-1} - h_na_{n-1,n} = x_{n-1} + h_{n-1}a_{n-1,n-1},
\]
wodurch $x_{n-1}$ und $h_{n-1}$ völlig bestimmt sind. Die nächst
vorhergehende Gleichung ergiebt
\[
y_{n-2} - h_na_{n-2,n} - h_{n-1}a_{n-2,n-1} = x_{n-2} + h_{n-2}a_{n-2,n-2},
\]
wodurch $x_{n-2}$ und $h_{n-2}$ völlig bestimmt sind, und so fahren wir
fort, bis alle $x_i$ und $h_i$ bestimmt sind.

Hiernach können wir zweckmässig die Gesammtheit der Zahlen $\xi$ oder
irgend ein System mit diesen congruenter Zahlen ein \gesperrt{volles
Restsystem nach dem Modul $\mu$ nennen.}

Die Anzahl der Individuen eines vollen Restsystems für den Modul $\mu$
ist gleich der absoluten Norm von $\mu$.

\begin{center}
\S.\,149.

\gesperrt{Congruenzen.}
\end{center}

Aus der Existenz eines vollen Restsystems nach einem Modul $\mu$, die
wir im vorigen Paragraphen nachgewiesen haben, ergiebt sich eine Reihe
von Sätzen, die mit bekannten elementaren Sätzen der rationalen
Zahlentheorie analog sind.

Bedeutet $\alpha$ irgend eine ganze Zahl in $\Omega$ und $\mu$ ein als
Modul dienendes ganzes Functional, und ist $\alpha$ relativ prim zu
$\mu$, so sind zwei Zahlen $\alpha\xi$ und $\alpha\xi'$ nur dann
congruent nach dem
%CONT>>
